%% sec-1-intro.tex -- Introduction

\section{Introduction}
\label{sec:intro}

Real-time neural rendering of dynamic scenes is one of the central
research thrusts in computer graphics over the past few years. Since
its introduction in 2023, 3D Gaussian Splatting (3DGS)~\cite{kerbl20233dgs}
has demonstrated that an explicit point-cloud representation combined
with differentiable rasterization can deliver 1080p real-time radiance
field rendering on desktop GPUs. Extending 3DGS to dynamic scenes
(4DGS), however, raises three core challenges: (\emph{i})~the choice of
temporal parameterization; (\emph{ii})~the storage cost of separating
dynamic and static content; and (\emph{iii})~the compute, bandwidth,
and memory budget of mobile GPUs.

This survey focuses on \emph{real-time 4DGS on mobile devices} and
systematizes more than 50 representative works published between 2023
and the first half of 2026. We organize the field around the following
guiding questions:
\begin{itemize}
  \item What 4DGS representation routes exist, and what are their
        trade-offs?
  \item How can the training stage (pruning, quantization, temporal
        redundancy) be compressed?
  \item How can the rendering stage (rasterization, ray tracing, mesh
        extraction) be accelerated?
  \item What are the current bottlenecks of Snapdragon / Adreno, and
        what are the most promising paths forward?
\end{itemize}

\noindent\textbf{Relationship to the project README.} This survey
complements \texttt{README.md}---the survey emphasizes academic depth
across the research landscape, while the README documents the
engineering path toward a production-ready mobile 4DGS renderer on
Snapdragon 8 Gen 4 / Adreno 830 + Vulkan 1.3. The survey targets
researchers; the README targets engineers.

\bigskip
\begin{itemize}
\item \textbf{2025-youn-success-gs}~\cite{youn2025successgs}: The two
      largest practical bottlenecks of 3DGS / 4DGS are storage cost
      (millions of Gaussians) and per-Gaussian motion inference
      overhead.
\end{itemize}
\noindent\textbf{Index of paper notes backing this section}:
\texttt{INDEX.md} \S A.

\bigskip
\noindent\textbf{Research-question roadmap.} The four guiding questions
are answered in the following sections:
\begin{center}
\begin{tabular}{ll}
\toprule
\textbf{Question} & \textbf{Answered in} \\
\midrule
Q1: 4DGS representation routes and trade-offs? & \S\ref{sec:representation} \\
Q2: Training-time compression routes? & \S\ref{sec:training} \\
Q3: Rendering-time acceleration routes? & \S\ref{sec:rendering} \\
Q4: Mobile deployment + future directions? & \S\ref{sec:mobile}, \S\ref{sec:datasets}, \S\ref{sec:discussion} \\
\bottomrule
\end{tabular}
\end{center}

\begin{figure}[t]
\centering
\small
\begin{tabular}{p{0.18\linewidth}p{0.72\linewidth}}
\toprule
\textbf{Faction} & \textbf{Scope and representative work} \\
\midrule
\textbf{A. 4DGS representation} &
Canonical + deformation (\cite{wu20244dgs}, \cite{yang2023deformable3dgs}); dynamic / static separation (\cite{yuan20254dgs1k}, \cite{zhang2024mega4dgsacceleration}); continuous-time (\cite{wang2026retimegs}); memory-efficient (\cite{shi2025sparse4dgs}, \cite{ren2026cubifygs}). \\
\textbf{B. 4DGS training accel.} &
Temporal pruning (\cite{lee2025omg4}, \cite{tu2025speede3dgs}); contextual / progressive coding (\cite{chen20254dgscc}, \cite{li2026pd4dgs}); streamable layered coding (\cite{yin2026cags}, \cite{zheng20254dgcpro}). \\
\textbf{C. 3DGS rendering accel.} &
Mobile rasterisation (\cite{du2026mobilegs}, \cite{du2026fluxgs}, \cite{huang2025seele}); ray tracing (\cite{liu20254dgrt}, \cite{ghosh2026gsnfs}); mesh / mobile neural rendering (\cite{feng2025lumina}); hardware acceleration (\cite{oh2025neo}). \\
\textbf{D. Mobile + streaming} &
Streaming codecs (\cite{li2026pd4dgs}, \cite{ke2025streamstgs}, \cite{veicht2026zipsplat}, \cite{yu2026codecsplat}); on-device decode (\cite{ghosh2026gsnfs}, \cite{shi2026evogs}); LoD (\cite{shi2026evogs}). \\
\textbf{E. Static 3DGS / cross-disciplinary} &
Compression primers (\cite{fan2023lightgaussian}, \cite{chen2024hacpp}, \cite{navaneet2023compact3d}); memory-efficient inference (\cite{liu2024efficientgs}, \cite{li2026vedal}); cross-domain (\cite{huang2026gaussianfluent}, \cite{feng2025lumina}). \\
\bottomrule
\end{tabular}
\caption{Faction taxonomy used throughout this survey. Factions A--D
represent the four primary dimensions of mobile 4DGS research; faction
E collects orthogonal static-3DGS acceleration and cross-disciplinary
work that informs the mobile 4DGS design space. The body of the survey
organises works into 9 sections (\S2--\S9) following the
``representation / training / rendering / deployment''
pipeline-stage axis, with cross-references back to the appropriate
faction where needed.}
\label{fig:taxonomy}
\end{figure}
